\section{Rationality and Stability in Classical Economics}
\label{sec:classicalrationality}
%%%%%%%%%%%%%%%%%%%%%%%%%%%%%%%%%%%%%%%%%%%%%%%%%%%%%%%%%%%%%%%%%%%%%%%%%%%%%%%

Economic rationality has been defined, refined, and contested for over two centuries.
This section traces the evolution of the concept from classical political economy
through neoclassical optimization to modern equilibrium theory,
identifying both the achievements and the limitations that motivate our framework.

\subsection{The Classical Foundations: Smith to Mill}

\paragraph{Adam Smith's Dual Insight.}
The \emph{Wealth of Nations} \citep{smith1776} established two foundational ideas:
\begin{enumerate}
  \item \textbf{Self-interest as mechanism:} ``It is not from the benevolence of the butcher, 
        the brewer, or the baker that we expect our dinner, but from their regard to their own interest.''
  \item \textbf{The invisible hand:} Individual self-interest, channeled through markets,
        produces socially beneficial outcomes.
\end{enumerate}

What is often overlooked is Smith's complementary insight in the \emph{Theory of Moral Sentiments} \citep{smith1759}:
human beings are not merely self-interested but also possess ``sympathy''---
the capacity to share others' feelings and to care about their judgment.
Smith's agent is already \emph{interdependent}: utility depends on others' states and opinions.

\paragraph{Example: The Butcher's Dilemma.}
Consider Smith's butcher in a small village:
\begin{itemize}
  \item Pure self-interest: Maximize profit by selling low-quality meat at high prices
  \item With reputation: Quality matters because repeat customers depend on trust
  \item With sympathy: The butcher cares about customers' welfare, not just their money
\end{itemize}

The classical tradition recognized that stable commerce requires more than self-interest---
it requires \emph{trust}, which emerges from repeated interaction and shared norms.
In our framework: stable exchange requires $C_{ij} > 0$ (complementary interests)
and context $\Psi$ that supports cooperation.

\paragraph{John Stuart Mill's Refinement.}
Mill \citep{mill1848} distinguished between ``higher'' and ``lower'' pleasures,
introducing a qualitative dimension to utility.
This anticipates the FEPSDE distinction:
financial gain (lower) vs. emotional fulfillment or identity satisfaction (higher).

Mill also emphasized that preferences are \emph{formed} through education and experience---
they are not exogenous but shaped by context.
This is precisely the endogeneity of $\Psi$ in our framework.

\subsection{The Marginalist Revolution: From Value to Utility}

\paragraph{The 1870s Breakthrough.}
Jevons, Menger, and Walras independently developed marginal utility theory,
shifting economics from objective labor value to subjective utility.

Key innovations:
\begin{itemize}
  \item Utility is measurable (at least ordinally)
  \item Rationality means maximizing utility subject to constraints
  \item Equilibrium is achieved when marginal utilities equal price ratios
\end{itemize}

\paragraph{Example: The Diamond-Water Paradox Resolved.}
Why are diamonds expensive and water cheap, despite water being essential?

Classical answer (labor theory): Diamonds require more labor to produce.

Marginalist answer: Price reflects \emph{marginal} utility, not total utility.
Water is abundant; the marginal glass adds little. Diamonds are scarce; the marginal diamond adds much.

This resolution required the concept of diminishing marginal utility \citep{marshall1890}:
$\partial^2 U / \partial x^2 < 0$.
In our notation (see Appendix~G for full glossary): the \emph{own-complementarity} $C_{ii} < 0$ (concavity).

\paragraph{What the Marginalists Assumed Away.}
The marginalist framework achieved elegance by assuming:
\begin{itemize}
  \item \textbf{Separability:} $\partial^2 U_i / \partial x_i \partial x_j = 0$ (no cross-effects)
  \item \textbf{Stability:} Preferences don't change
  \item \textbf{Independence:} My utility doesn't depend on your consumption
\end{itemize}

These assumptions enabled mathematical tractability but excluded precisely
the interdependencies that drive real economic behavior.

\subsection{General Equilibrium: The Arrow-Debreu Achievement}

\paragraph{The Existence Theorem.}
\citet{arrowdeb} proved---and \citet{debreu1959} formalized axiomatically---that under certain conditions,
a competitive equilibrium exists and is Pareto efficient.

Required conditions:
\begin{enumerate}
  \item Convex preferences (diminishing marginal utility)
  \item Convex production sets (diminishing returns)
  \item Complete markets (all goods can be traded)
  \item Perfect information (all agents know all prices)
  \item No externalities (no unpriced interdependencies)
\end{enumerate}

\paragraph{Example: The Edgeworth Box.}
Two agents, two goods. Each starts with an endowment.
Trade moves them to a Pareto-efficient allocation on the contract curve.
Competitive prices ensure market clearing.

This is the paradigm of classical rationality:
given preferences and endowments, the invisible hand finds the efficient allocation.

\paragraph{What Arrow-Debreu Achieved.}
\begin{itemize}
  \item Rigorous proof that markets \emph{can} work
  \item Clear conditions for efficiency
  \item Foundation for welfare economics
  \item Benchmark for policy analysis
\end{itemize}

\paragraph{What Arrow-Debreu Assumed.}
The fifth condition---no externalities---is crucial.
It means:
\begin{equation}
  C_{ij} = \frac{\partial^2 U_i}{\partial x_i \partial x_j} = 0 
  \quad \text{for all } i \neq j
\end{equation}

Your consumption doesn't affect my utility.
Your production doesn't affect my costs.
We are connected only through prices.

This is the \emph{separability assumption} that our framework relaxes.

\subsection{Stability: The Missing Dimension}

\paragraph{Efficiency vs. Stability.}
Arrow-Debreu proves that equilibrium is efficient.
But is it stable? Will the economy return to equilibrium after a shock?

\paragraph{The Sonnenschein-Mantel-Debreu Problem.}
\citet{sonnenschein1972} and subsequent work showed that
aggregate excess demand functions can have almost any shape
consistent with Walras' law and homogeneity.

Implication: General equilibrium theory places almost no restrictions on dynamics.
Multiple equilibria are possible. Stability is not guaranteed.

\paragraph{Example: Coordination Failure.}
Consider two workers choosing whether to specialize.
If both specialize, they can trade and both gain.
If one specializes and the other doesn't, the specialist loses.

\begin{center}
\renewcommand{\arraystretch}{1.2}
\begin{tabular}{l|cc}
 & Worker 2: Specialize & Worker 2: Autarky \\
\hline
Worker 1: Specialize & (3, 3) & (0, 2) \\
Worker 1: Autarky & (2, 0) & (2, 2) \\
\end{tabular}
\end{center}

Two equilibria: (Specialize, Specialize) and (Autarky, Autarky).
The first is Pareto-superior, but the second is also stable.
Arrow-Debreu cannot predict which will emerge.

In our framework: Both have $K \approx 1$ (coherent), but different $Q$.
The economy can be stably trapped in the inferior equilibrium.

\paragraph{The Stability Gap.}
Classical rationality defines an end-state property---a limitation Keynes \citep{keynes1936} emphasized in challenging the assumption of automatic market clearing:
optimal allocation given preferences and constraints.
But it says nothing about:
\begin{itemize}
  \item How equilibrium is reached
  \item Whether equilibrium persists under perturbation
  \item How the economy moves between multiple equilibria
  \item What happens when preferences themselves change
\end{itemize}

This is the gap our framework addresses.

\subsection{Evolutionary Game Theory: Selection Without Rationality}

\paragraph{The Intuition.}
Classical economics assumes people \emph{calculate} the optimal choice. 
But nobody solves differential equations while shopping. 
Nobody knows all alternatives. Nobody has infinite computing power.

So where does ``rational'' behavior come from?

The evolutionary answer: \textbf{selection, not calculation}.
You don't need to know what's optimal. You just need to survive.
Those who play bad strategies disappear. Those who play good ones spread.
Rationality emerges as a \emph{result}, not as an assumption.

Think of it this way: A chess grandmaster doesn't calculate every possible move.
She uses patterns learned over years of play---patterns that survived because they won games.
The market works similarly: firms that make bad choices go bankrupt; 
strategies that work get imitated and spread.

\paragraph{Why This Matters for the Framework.}
This shift has profound implications:
\begin{itemize}
  \item \textbf{Bounded rationality is not a bug:} Heuristics and biases evolved 
        because they worked in ancestral environments. $\Psi_C$ reflects evolutionary
        calibration, not ``irrationality.''
  \item \textbf{Social preferences can survive:} Altruism seems costly, but evolution 
        found ways to sustain it (kin selection, reciprocity, group selection).
        $\Psi_S$ is evolutionarily stable.
  \item \textbf{Multiple equilibria are natural:} Selection can stabilize different 
        configurations. History determines which one emerges.
\end{itemize}

\paragraph{The Technical Foundation.}
The stability gap motivated a fundamentally different approach: rather than assuming
agents calculate optimal strategies, evolutionary game theory asks which strategies
\emph{survive} competitive selection \citep{weibull1995,samuelson1997,gintis2000}.

\paragraph{Maynard Smith: Evolutionary Stable Strategies.}
\citet{maynardsmith1982} introduced the concept of an Evolutionarily Stable Strategy (ESS):
a strategy that, once prevalent, cannot be invaded by any rare mutant strategy.

\begin{definition}[ESS]
Strategy $s^*$ is evolutionarily stable if for all $s \neq s^*$:
\begin{equation}
  \pi(s^*, s^*) > \pi(s, s^*) \quad \text{or} \quad
  [\pi(s^*, s^*) = \pi(s, s^*) \text{ and } \pi(s^*, s) > \pi(s, s)]
\end{equation}
where $\pi(s_i, s_j)$ is the payoff to strategy $s_i$ against $s_j$.
\end{definition}

\paragraph{Key Insight.}
ESS refines Nash equilibrium: every ESS is a Nash equilibrium, but not vice versa.
ESS adds a \emph{stability} requirement---the equilibrium must be robust to invasion.

\paragraph{Replicator Dynamics.}
The standard model of evolutionary dynamics \citep{taylor1978}:
\begin{equation}
  \dot{x}_i = x_i \left[ \pi_i(x) - \bar{\pi}(x) \right]
  \label{eq:replicator}
\end{equation}
where $x_i$ is the population share using strategy $i$, $\pi_i$ is its payoff,
and $\bar{\pi}$ is the population average payoff.

Strategies that earn above-average payoffs grow; those below average shrink.
This is selection without foresight---agents don't optimize, they are selected.

\paragraph{Connection to Classical Equilibrium.}
\citet{weibull1995} showed:
\begin{itemize}
  \item Rest points of replicator dynamics are Nash equilibria
  \item Stable rest points are refinements of Nash (often ESS)
  \item Dynamics provide equilibrium selection among multiple Nash equilibria
\end{itemize}

\paragraph{Bounded Rationality as Evolutionary Success.}
\citet{gintis2009} argued that bounded rationality can be understood evolutionarily:
heuristics and biases that seem ``irrational'' may be fitness-maximizing in ancestral environments.
This connects $\Psi_C$ (cognitive context) to evolutionary foundations.

\paragraph{Social Preferences and Evolution.}
How can costly prosocial behavior survive selection? Several mechanisms:
\begin{itemize}
  \item \textbf{Kin selection:} Altruism toward relatives \citep{hamilton1964}
  \item \textbf{Reciprocity:} Conditional cooperation in repeated games \citep{axelrod1984,nowak2006}
  \item \textbf{Group selection:} Groups with cooperators outcompete groups without \citep{boydricherson2005}
  \item \textbf{Indirect reciprocity:} Reputation sustains cooperation \citep{nowak1998}
\end{itemize}

\citet{fehr2004nature} showed that altruistic punishment---costly punishment of defectors---can
stabilize cooperation even in one-shot anonymous interactions.

\paragraph{Quantal Response Equilibrium.}
\citet{mckelvey1995} bridged rationality and bounded rationality with QRE:
\begin{equation}
  P(s_i) = \frac{\exp(\lambda \cdot \pi_i)}{\sum_j \exp(\lambda \cdot \pi_j)}
  \label{eq:qre}
\end{equation}
where $\lambda$ measures rationality: $\lambda = 0$ is random, $\lambda \to \infty$ is fully rational.

QRE explains systematic deviations from Nash in experiments while maintaining
equilibrium logic. It operationalizes $\Psi_C$: lower cognitive capacity means lower $\lambda$.

\paragraph{Framework Integration.}
Evolutionary game theory provides the dynamic foundation for our framework:
\begin{itemize}
  \item $C^*$ is not calculated but \emph{emerges} from evolutionary dynamics
  \item Multiple $C^*$ may exist; selection determines which is reached
  \item Bounded rationality ($\Psi_C < \infty$) is evolutionarily grounded
  \item Social preferences ($\Psi_S$) are evolutionarily stable under specific conditions
\end{itemize}

The equilibrium concept shifts from ``what rational agents would choose'' to
``what survives selection''---a more robust foundation that doesn't require
implausible computational abilities.

\subsection{From Efficiency to Coherence}

\paragraph{The Limits of Optimization.}
Classical rationality is agent-centric: each agent maximizes given the environment.
But when agents' utilities are \emph{interdependent}, optimization by each
doesn't guarantee system-level stability.

\paragraph{Example: The Tragedy of the Commons.}
Each herder rationally adds cattle to the common pasture.
Each optimizes individually.
The pasture is destroyed---a collectively irrational outcome.

Arrow-Debreu would say: ``This is an externality; define property rights.''
But property rights themselves require institutions, which require context $\Psi$.

\paragraph{Example: Bank Runs.}
Each depositor rationally withdraws funds when others do \citep{diamonddybvig1983}.
Individual optimization leads to collective collapse.
The efficient equilibrium (all deposit) is unstable.

In our framework: The ``all deposit'' equilibrium has high $Q$ but low $K$
when trust is fragile. A small shock triggers a cascade to the ``all withdraw'' equilibrium.

\paragraph{The Coherence Criterion.}
We propose that rationality should refer not only to individual optimization
but to the \emph{stability of the system of interdependent utilities}:
\begin{itemize}
  \item How robust is the equilibrium to shocks?
  \item How well-aligned are agents' expectations with realized outcomes?
  \item How consistent are institutions, norms, and behavior?
\end{itemize}

This is what the coherence index $K$ measures.

\subsection{Historical Examples of Coherence and Incoherence}

\paragraph{The Gold Standard (1870-1914): High K, Moderate Q.}
The classical gold standard was remarkably coherent:
\begin{itemize}
  \item Clear rules: currencies pegged to gold
  \item Automatic adjustment: gold flows balanced payments
  \item Shared expectations: credible commitment
\end{itemize}

$K \approx 1$: The system was internally consistent.
$Q$ was moderate: growth was constrained by gold supply; crises were deflationary.

\paragraph{Weimar Germany (1921-1923): Collapse of K.}
Hyperinflation destroyed coherence:
\begin{itemize}
  \item Monetary rules: broken by war finance
  \item Expectations: deanchored, self-fulfilling collapse
  \item Institutions: lost credibility
\end{itemize}

$K \to 0$: Complete breakdown of the institutional framework.
$Q$ collapsed: economic and social devastation.

\paragraph{Post-War Reconstruction (1945-1970): High K, High Q.}
The Bretton Woods system and welfare state created:
\begin{itemize}
  \item Clear rules: fixed exchange rates, capital controls
  \item Shared expectations: growth, stability, inclusion
  \item Complementary institutions: monetary policy + fiscal policy + social insurance
\end{itemize}

$K \approx 1$: Coherent configuration of all elements.
$Q$ high: ``Les Trente Glorieuses,'' unprecedented prosperity.

\paragraph{The 2008 Financial Crisis: K Collapse.}
Before the crisis:
\begin{itemize}
  \item Apparent coherence: Rising asset prices, low volatility
  \item Hidden incoherence: Risk models assumed stable correlations
\end{itemize}

During the crisis:
\begin{itemize}
  \item Correlation breakdown: ``All assets fell together''
  \item Expectation collapse: Counterparty trust evaporated
  \item Institutional failure: Models proved wrong
\end{itemize}

$K_{market}$: Dropped from $\approx 0.95$ to $\approx 0.3$.
Recovery required massive intervention to restore coherence.

\subsection{Why Large Language Models Reproduce Classical Predictions}

\paragraph{The Training Corpus Problem.}
The limitations of classical rationality help explain a striking empirical finding:
when large language models (LLMs) are used to simulate human behavior, they often
reproduce \emph{classical} predictions rather than actual human behavior.

\citet{peng2025} found that LLM-based ``digital twins'' replicated only 50\% of
classic behavioral economics experiments. More tellingly, the failures were
systematic: LLMs often predicted the \emph{textbook rational} response when
humans showed behavioral deviations.

\paragraph{Example: The Default Effect.}
In organ donation experiments, behavioral economics predicts that defaults
strongly influence choices---people tend to accept whatever is pre-selected
\citep{johnson2003}. But in Peng et al.'s replication:
\begin{itemize}
    \item \textbf{Humans:} Showed no significant default effect in this context
    \item \textbf{Digital twins:} Showed a strong default effect
\end{itemize}

The LLM had ``learned'' from its training corpus that defaults matter---this is
propositional knowledge accurately extracted from behavioral economics literature.
But it could not predict \emph{when} defaults matter (context-dependence on $\Psi$).

\paragraph{Why This Happens.}
LLMs are trained on text, including economics textbooks, Wikipedia articles, and
academic papers. This corpus is dominated by:
\begin{enumerate}
    \item Classical theory: Presented as the benchmark
    \item Behavioral findings: Presented as ``deviations'' from the benchmark
    \item Stylized descriptions: ``The endowment effect is robust''
\end{enumerate}

What is \emph{not} in the training corpus:
\begin{itemize}
    \item Quantitative effect sizes as functions of context
    \item The specific conditions under which effects appear or disappear
    \item The interaction structure between contextual factors
\end{itemize}

The LLM knows \emph{that} behavioral effects exist but not \emph{when} or \emph{how strongly}.
When forced to make a prediction, it often defaults to the classical benchmark
that dominates its training data.

\paragraph{The Irony.}
This creates a profound irony: LLMs trained on behavioral economics literature
often reproduce the rational-choice predictions that behavioral economics was
invented to correct. The models have learned the \emph{words} but not the
\emph{phenomena}.

\paragraph{Implications for This Framework.}
This finding motivates the empirical calibration approach developed in Section~4.X.
Rather than deriving behavior from axioms (which LLMs can reproduce) or describing
effects in words (which LLMs can parrot), we calibrate on the \emph{data themselves}---
the actual variation in behavior across contexts that reveals the functional form
of $C^*(\Psi)$.

The classical tradition achieved rigor through axiomatization.
The behavioral tradition documented deviations through experiments.
This framework seeks to combine both: rigorous predictions calibrated on experimental evidence.

\subsection{Summary: What Classical Economics Gets Right and Wrong}

\paragraph{Achievements of Classical Rationality.}
\begin{itemize}
  \item Rigorous framework for optimization
  \item Proof that markets can achieve efficiency
  \item Clear welfare criteria
  \item Foundation for policy analysis
\end{itemize}

\paragraph{Limitations of Classical Rationality.}
\begin{itemize}
  \item Assumes away interdependence ($C_{ij} = 0$)
  \item Treats preferences as exogenous (no $\Psi$ dynamics)
  \item Describes efficiency but not stability
  \item Cannot explain multiple equilibria selection
  \item Silent on institutional design
\end{itemize}

\paragraph{The Extension Required.}
A complete theory of rationality must include:
\begin{enumerate}
  \item Utility interdependence: $C_{ij} \neq 0$
  \item Endogenous context: $\Psi_{t+1} = f(\Psi_t, a_t)$
  \item Stability criterion: $K$ as measure of coherence
  \item Quality criterion: $Q$ as measure of welfare
  \item Multi-level analysis: Individual to global
\end{enumerate}

This is what the $(C, \Psi)$ framework provides.

\paragraph{From Theory to Calibration.}
Classical economics derived predictions from axioms.
Behavioral economics documented deviations through experiments.
This framework takes the next step: calibrating $C^*(\Psi)$ on the accumulated
experimental record to generate \emph{quantitative} predictions about \emph{when}
classical assumptions hold and when they fail.

Section~3 examines the specific limits of utility maximization that motivated
the behavioral revolution. Section~4 develops the empirical foundations for
measuring context and calibrating the framework---including Section~4.X, which
explains why this calibration approach succeeds where LLM-based simulation fails.

%%%%%%%%%%%%%%%%%%%%%%%%%%%%%%%%%%%%%%%%%%%%%%%%%%%%%%%%%%%%%%%%%%%%%%%%%%%%%%%
